\documentclass[11pt]{article}

% cs250preamble.tex --- shared preamble for CS 250 course notes
%
% This file is \input by main.tex and by each individual module
% file so that each module can still compile standalone. Any
% package or custom-environment change should be made here, not
% in the individual module files.

\usepackage[T1]{fontenc}
\usepackage[margin=1in]{geometry}
\usepackage{amsmath, amssymb, amsthm}
\usepackage{enumitem}
\usepackage{hyperref}
\usepackage{booktabs}
\usepackage{listings}
\usepackage{xcolor}
\usepackage{tikz}
\usetikzlibrary{shapes.geometric, shapes.gates.logic.US, shapes.gates.logic.IEC, arrows.meta, positioning, calc, trees}
\usepackage{bussproofs}
\usepackage{mdframed}

\lstset{
  basicstyle=\ttfamily\small,
  frame=single,
  breaklines=true,
  columns=fullflexible,
  mathescape=true
}

\theoremstyle{definition}
\newtheorem{theorem}{Theorem}[section]
\newtheorem{definition}[theorem]{Definition}
\newtheorem{example}[theorem]{Example}

\hypersetup{
  colorlinks=true,
  linkcolor=blue!60!black,
  urlcolor=blue!60!black
}

% Key-takeaway box: used at the end of major sections and at the end of
% each module to summarise the main points. Expects \item bullets inside.
\newmdenv[
  linewidth=0.8pt,
  linecolor=blue!40!black,
  backgroundcolor=blue!3,
  skipabove=8pt,
  skipbelow=8pt,
  innertopmargin=7pt,
  innerbottommargin=7pt,
  innerleftmargin=9pt,
  innerrightmargin=9pt
]{keytakeawaybox}
\newenvironment{keytakeaway}{%
  \begin{keytakeawaybox}%
  \noindent\textbf{Key takeaways.}%
  \begin{itemize}[leftmargin=1.3em, topsep=2pt, itemsep=1pt, label=\textbullet]%
}{%
  \end{itemize}%
  \end{keytakeawaybox}%
}

% Summary/looking-ahead environment: used once at the end of each module
% to wrap up what the module built and point forward.
\newmdenv[
  linewidth=0.8pt,
  linecolor=gray!60,
  backgroundcolor=gray!5,
  skipabove=8pt,
  skipbelow=8pt,
  innertopmargin=7pt,
  innerbottommargin=7pt,
  innerleftmargin=9pt,
  innerrightmargin=9pt
]{summarybox}

% Sidebar environment: used for tangential asides---historical notes,
% pointers to other areas of math/CS, cross-paradigm remarks---that
% enrich the main narrative without being essential to it. Takes one
% argument: the sidebar's title.
\newmdenv[
  linewidth=0.6pt,
  linecolor=gray!60,
  backgroundcolor=gray!4,
  skipabove=8pt,
  skipbelow=8pt,
  innertopmargin=6pt,
  innerbottommargin=6pt,
  innerleftmargin=9pt,
  innerrightmargin=9pt
]{sidebarbox}
\newenvironment{sidebar}[1]{%
  \begin{sidebarbox}%
  \noindent\textbf{Sidebar: #1.}\par\medskip%
}{%
  \end{sidebarbox}%
}

% Reading-map environment: used at the top of each numbered module to
% indicate Epp coverage, prerequisites, relevant intermezzi, and
% suggested practice problems. Expects four \rmentry{label}{body}
% items inside (or arbitrary paragraphs).
\newmdenv[
  linewidth=0.8pt,
  linecolor=black!55,
  backgroundcolor=yellow!4,
  skipabove=10pt,
  skipbelow=14pt,
  innertopmargin=9pt,
  innerbottommargin=9pt,
  innerleftmargin=11pt,
  innerrightmargin=11pt
]{readingmapbox}
\newcommand{\rmentry}[2]{\par\noindent\textbf{#1}\hspace{0.4em}#2\par}
\newenvironment{readingmap}{%
  \begin{readingmapbox}%
  \noindent{\bfseries\large Reading map}%
  \vspace{2pt}%
  \setlength{\parskip}{4pt plus 1pt}%
}{%
  \end{readingmapbox}%
}


\title{CS 250 --- Notation Summary}
\author{}
\date{}

\begin{document}
\maketitle

{\small
\begin{center}
\begin{tabular}{lll}
\toprule
\textbf{Symbol} & \textbf{Name} & \textbf{Introduced in} \\
\midrule
\multicolumn{3}{l}{\textit{Propositional Logic}} \\[2pt]
$\wedge$ & Conjunction (and) & Module 2 \\
$\vee$ & Disjunction (or) & Module 2 \\
$\lnot$ or $\sim$ & Negation (not) & Module 2 \\
$\to$ & Implication (if\ldots then) & Module 2 \\
$\leftrightarrow$ & Biconditional (if and only if) & Module 2 \\
$\equiv$ & Logical equivalence & Module 2 \\
$\oplus$ & Exclusive or (XOR) & Module 2 \\
T, F & Truth values & Module 2 \\
$\bot$ & Absurdity / falsum & Gentzen Intermezzo \\
\midrule
\multicolumn{3}{l}{\textit{Quantifiers and Predicates}} \\[2pt]
$\forall$ & For all (universal quantifier) & Module 4 \\
$\exists$ & There exists (existential quantifier) & Module 4 \\
$P(x)$ & Predicate applied to $x$ & Module 4 \\
\midrule
\multicolumn{3}{l}{\textit{Sets}} \\[2pt]
$\{\ \}$ & Set braces & Module 1 \\
$\in$, $\notin$ & Element of, not element of & Module 1 \\
$\subseteq$ & Subset & Module 1 \\
$\cap$ & Intersection & Module 1 \\
$\cup$ & Union & Module 1 \\
$\setminus$ & Set difference & Module 1 \\
$\overline{A}$ & Complement & Module 1 \\
$\emptyset$ & Empty set & Module 1 \\
$|A|$ & Cardinality & Module 1 \\
$\mathcal{P}(A)$ & Power set & Module 1 \\
$A \times B$ & Cartesian product & Module 1 \\
$B^A$ & Function space ($A \to B$) & Cardinality Intermezzo \\
$\{x \mid P(x)\}$ & Set-builder notation & Module 1 \\
\midrule
\multicolumn{3}{l}{\textit{Functions and Relations}} \\[2pt]
$f : A \to B$ & Function from $A$ to $B$ & Module 1 \\
$a\, R\, b$ & $a$ is related to $b$ by $R$ & Module 1 \\
$\cong$ & Isomorphism / bijection exists & Cardinality Intermezzo \\
\bottomrule
\end{tabular}
\end{center}
}

\clearpage

{\small
\begin{center}
\begin{tabular}{lll}
\toprule
\textbf{Symbol} & \textbf{Name} & \textbf{Introduced in} \\
\midrule
\multicolumn{3}{l}{\textit{Number Theory}} \\[2pt]
$\mathbb{N}$ & Natural numbers & Module 1 \\
$\mathbb{Z}$ & Integers & Module 5 \\
$\mathbb{Q}$ & Rationals & Module 5 \\
$\mathbb{R}$ & Reals & Module 4 \\
$a \mid b$ & $a$ divides $b$ & Module 5 \\
$\gcd(a,b)$ & Greatest common divisor & Module 7 \\
$\mathbb{Z}/n\mathbb{Z}$ & Integers mod $n$ & Module 6 \\
$\lfloor x \rfloor$, $\lceil x \rceil$ & Floor, ceiling & Module 6 \\
\midrule
\multicolumn{3}{l}{\textit{Sequences and Sums}} \\[2pt]
$a_n$ & $n$th term of a sequence & Module 8 \\
$\sum$ & Summation & Module 8 \\
$\prod$ & Product & Module 8 \\
$n!$ & Factorial & Module 8 \\
\midrule
\multicolumn{3}{l}{\textit{Data Structures}} \\[2pt]
Nil, Cons & List constructors & Module 10 \\
Leaf, Node & Binary tree constructors & Module 10 \\
\midrule
\multicolumn{3}{l}{\textit{Proof Notation}} \\[2pt]
$\therefore$ & Therefore & Module 3 \\
$\square$ or QED & End of proof & Throughout \\
$\{P\}\,C\,\{Q\}$ & Hoare triple & Module 7 \\
\bottomrule
\end{tabular}
\end{center}
}

\end{document}
